% Generated draft: requires mathematical and source review.
\documentclass[11pt]{article}
\usepackage[margin=25mm]{geometry}
\usepackage{fontspec}
\setmainfont{Latin Modern Roman}
\usepackage{amsmath,amssymb,mathtools,unicode-math}
\setmathfont{Latin Modern Math}
\usepackage{enumitem,needspace,xurl,hyperref}
\setlength{\parindent}{0pt}
\setlength{\parskip}{5pt}
\setlength{\emergencystretch}{3em}
\newcommand{\N}{\mathbb{N}}
\newcommand{\Z}{\mathbb{Z}}
\newcommand{\Q}{\mathbb{Q}}
\newcommand{\R}{\mathbb{R}}
\newcommand{\C}{\mathbb{C}}
\newcommand{\F}{\mathbb{F}}
\newcommand{\A}{\mathbb{A}}
\newcommand{\PP}{\mathbb{P}}
\newcommand{\E}{\mathbb{E}}
\newcommand{\eps}{\varepsilon}
\newcommand{\dd}{\,\mathrm{d}}
\begin{document}
\section*{UnsolvedMath ID: 30005119}
\textbf{Source problem code:} catalogue entry
\subsection{Definitions and mathematical statement}
Let \(k\geq 6\) be even, let \(q\) be a prime power, and let \(a=(a_1,\ldots,a_k)\) and \(b=(b_1,\ldots,b_k)\) belong to \(\F_q^k\). Put
\[
L=\begin{pmatrix}a_1&\cdots&a_k\\b_1&\cdots&b_k\end{pmatrix}.
\]
The hypothesis on the scalar solution space is
\[
\dim_{\F_q}\ker L=k-2,
\]
equivalently, \(\operatorname{rank}_{\F_q}L=2\).
For each positive integer \(n\), define
\[
S_L(n)=\left\{(x_1,\ldots,x_k)\in(\F_q^n)^k:
\sum_{i=1}^k a_i x_i=0,\quad \sum_{i=1}^k b_i x_i=0\right\}.
\]
For \(f:\F_q^n\to[0,1]\), let
\[
M_L(f)=\frac{1}{|S_L(n)|}\sum_{(x_1,\ldots,x_k)\in S_L(n)}
\left(\prod_{i=1}^k f(x_i)+\prod_{i=1}^k(1-f(x_i))\right).
\]
Thus the average is uniform over all solutions, with repeated coordinates permitted. Following the supplied definition, the pair is \emph{common with respect to two colours over \(\F_q\)} if
\[
\forall n\geq1\;\forall f:\F_q^n\to[0,1],\qquad M_L(f)\geq 2^{1-k}.
\]
It is \emph{uncommon} if
\[
\exists n\geq1\;\exists f:\F_q^n\to[0,1],\qquad M_L(f)<2^{1-k}.
\]
\textbf{Source question (Conjecture 16).} For every even \(k\geq6\), does the rank-two hypothesis imply that the pair is uncommon over every sufficiently large finite field? For a fixed field and coefficient pair, the proposed conclusion is precisely the last displayed strict inequality for some \(n\) and \(f\).

\textbf{Unresolved formulation qualifications.} The record does not specify how the coefficients are fixed or transported as the field varies, or whether the threshold for \(q\) may depend on the coefficient pair as well as on \(k\). Consequently, a unique quantifier order for the large-field assertion cannot be recovered. The supplied formulation imposes only the rank condition: it gives no additional restrictions on zero coefficients, equation supports, or linear combinations of the two equations. No such restrictions are added here; whether further hypotheses were intended requires mathematical review.
\subsection{Short English statement}
The supplied conjecture asks whether every pair of independent homogeneous linear equations in an even number \(k\geq6\) of variables is uncommon over sufficiently large finite fields. Uncommonness means that some dimension and some fractional two-colouring give a monochromatic solution average strictly below \(2^{1-k}\). The coefficient conventions, dependence of the field-size threshold, and possible missing hypotheses remain unresolved.
\subsection{Sources}
\begin{itemize}
\item Supplied local UnsolvedMath registry record, permanent ID \(30005119\), title \emph{Uncommon Pairs of Linear Equations}, exact source code \texttt{catalogue entry}; collection: Oberwolfach Reports Open Problems. The registry records proposed year \(2022\).
\item Supplied report citation: \emph{Combinatorics, Probability and Computing} (2023), Oberwolfach Reports, DOI \texttt{10.4171/owr/2022/22}, pp.\ 1165--1237, Conjecture 16. \url{https://doi.org/10.4171/owr/2022/22}. The extraction places the name Natasha Morrison immediately after the conjecture; the registry's proposer field is empty, so the precise attribution remains uncertain.
\item Altman and Liebenau, supplied reference label: ``minimal rank-two systems.'' \url{https://doi.org/10.1112/blms.70135}.
\item Dong, Li, and Zhao, supplied reference label: ``uncommon linear systems.'' \url{https://arxiv.org/abs/2404.17005}.
\item Supplied literature assessment, checked \texttt{2026-08-22}: status labelled partially solved. The curation reports uncommonness results for several broad subclasses, including generic systems and systems of girth \(k-1\), while describing the full rank-two conjecture as unresolved. This is supplied dated curation, not an independently verified literature assessment.
\item The registry labels the statement \texttt{corrected\_verified} and reports that the official report was checked directly. These are supplied provenance claims; no source visit or independent verification is asserted in this draft.
\end{itemize}
\end{document}
